\section{Discussion}
\label{sec:discussion}

\textbf{Why are absolute numbers so close to random?} Object \emph{state} -- deformation,
contact, orientation change -- is underrepresented in web-scale VLM training data
relative to object \emph{identity} and \emph{position}. Models can usually describe which
objects are present and roughly where, but not reliably what has happened to them or
which phase of an action is depicted. Compounding the problem, the models rarely use
``Cannot be determined'' as an answer even when it is correct, defaulting instead to a
confident Yes or No (\cref{sec:results-lora}: BASE Qwen3-4B scores 0--15\% on the CnBD
variant). The 70\%-plus \texttt{task\_success} accuracy of the Qwen3 family is a partial
exception: binary success/failure framing may be closer to the kind of outcome
classification VLMs see at web scale.

\textbf{Scale helps, but is not enough.} Within-family scaling consistently improves
accuracy (Gemma3 E2B~$\to$~E4B~$\to$~31B), but even 31B/32B-class models plateau in the
high-40s\% on phase-level reasoning. Cross-family, InternVL3-14B is competitive with
Qwen3-8B despite the parameter gap, suggesting architecture and training data matter at
least as much as raw parameter count. More parameters alone are unlikely to close the
gap; robot-domain training data is needed.

\textbf{LoRA gains require decomposition before deployment decisions.} The 11pp headline
improvement could be read as ``fine-tuning works for robot state reasoning,'' but the
decomposition in \cref{sec:results-lora} shows that roughly 80\% of the gain is two
coarse heuristics. The OOD-scene detector is genuinely useful and generalizable, but it is
domain-mismatch detection, not phase or temporal reasoning. The label-position bias on
\texttt{phase\_success} is net positive under this dataset's label distribution, but a
regression for the most safety-critical use case: detecting when a robot is currently in
the middle of an action. The lesson is to always decompose fine-tuning gains by question
type and answer class before concluding that a model is ``better'' at a target skill.

\textbf{Multi-view consistency tracks confidence, not correctness.} Models are most
self-consistent when uniformly right (72.3\%) or uniformly wrong (66.4\%) across views,
and least consistent when only some views are correct (34.4\%). We read this as
view-consistency reflecting how unambiguous a scene is to the model, not whether its
answer is correct. This connects to the scene-difficulty finding (\cref{sec:results}):
scenes that drive accuracy down also destabilize cross-view agreement, suggesting both
effects share the same root cause -- visual ambiguity, not task semantics. For multi-view
fusion, this implies that simple voting across views may not help on exactly the hardest
scenes, where inconsistency is highest.
